\section*{第 \ref{ch:b3:x-ray-diffraction} 章——晶体学与 X 射线衍射}

\begin{solution}{exo:b3:x-ray-diffraction:1}
$d = a/\sqrt N$: $d_{111} = \qty{325.66}{pm}$, $d_{200} = \qty{282.03}{pm}$, $d_{220} =
\qty{199.42}{pm}$.
\end{solution}

\begin{solution}{exo:b3:x-ray-diffraction:2}
面心立方：最先出现的反射为 (111) 和 (200)。$\sin\theta = \lambda\sqrt N/2a$：$2\theta =
\ang{43.32}$ 和 \ang{50.45}。
\end{solution}

\begin{solution}{exo:b3:x-ray-diffraction:3}
P：全部七个。I（$h + k + l$ 为偶数）：110、200、211、220。F（指标全奇或全偶）：111、200、220。
\end{solution}

\begin{solution}{exo:b3:x-ray-diffraction:4}
$\mathbf a^* = \mathbf e_x/a$，$\mathbf b^* = \mathbf e_y/b$：矩形点阵，边长为 $1/a$
和 $1/b$，正点阵的长边在倒点阵中成为短边。
\end{solution}

\begin{solution}{exo:b3:x-ray-diffraction:5}
$\sin^2\theta$ 之比为 $1 : 2 : 3 : 4 : 5$，即 $N = 2, 4, 6, 8, 10$：体心立方，
反射为 (110)、(200)、(211)、(220)、(310)。由第一行的数据可得 $d_{110} = \lambda/2\sin(\ang{20.135}) =
\qty{223.77}{pm}$，$a = d\sqrt2 = \qty{316.5}{pm}$：钨。
\end{solution}

\begin{solution}{exo:b3:x-ray-diffraction:6}
$Z = \rho N_Aa^3/M = 1.99 \times 6.022\times10^{23} \times (6.2929\times10^{-8})^3/74.6 = 4.0$.
\end{solution}

\begin{solution}{exo:b3:x-ray-diffraction:7}
$\beta = \qty{0.40}{\degree} = \qty{6.98e-3}{rad}$, $\cos\theta = \cos\ang{15.85} = 0.962$:
$L = 0.9 \times 0.15406/(6.98\times10^{-3} \times 0.962) = \qty{21}{nm}$.
\end{solution}

\begin{solution}{exo:b3:x-ray-diffraction:8}
$F_{hkl} = f_{\ce{Cs+}} + f_{\ce{Cl-}}(-1)^{h+k+l}$: $F_{100} = 54 - 18 = 36$, $F_{110} = 72$:
强度之比为 1 : 4。体心位置上的离子与顶点上的离子不同，
因此点阵为简单点阵（体心点阵要求两点处的内容相同），
且 (100) 反射存在。
\end{solution}

\begin{solution}{exo:b3:x-ray-diffraction:9}
每个原子都有一个平移 $(\frac12,\frac12,0)$ 的副本，使 $F$ 乘以 $1 +
\eu^{\iu\pi(h+k)} = 1 + (-1)^{h+k}$，当 $h + k$ 为奇数时为零。
\end{solution}

\begin{solution}{exo:b3:x-ray-diffraction:10}
准晶具有长程有序而无周期性：它们没有
点阵，因此针对点阵的限制不适用于它们。
其尖锐的衍射斑点促使人们将晶体定义为
具有本质上离散衍射图样的任何固体，无论是否周期性。
\end{solution}

\begin{solution}{exo:b3:x-ray-diffraction:11}
$c$ 滑移面将 $(x, y, z)$ 变为 $(x, \frac12 - y, \frac12 + z)$。当 $k = 0$ 时，两者的
相位因子为 $\eu^{2\pi\iu(hx + lz)}$ 和 $\eu^{2\pi\iu(hx + lz + l/2)} = (-1)^l
\eu^{2\pi\iu(hx+lz)}$，二者之和在 $l$ 为奇数时为零。
\end{solution}

\begin{solution}{exo:b3:x-ray-diffraction:12}
反演、镜面或滑移面会使分子变为其镜像，因此
含有其一的晶体将同时包含两种对映异构体。纯的对映异构体
仅能结晶于 65 个由旋转、螺旋轴与平移构成的
空间群（佐恩克群）中，例如 \textit{P}$2_12_12_1$。由弗里德尔定律可知
一种对映异构体的图样与另一种相同；绝对
构型可由反常散射引起的微小偏差求得，
这些偏差来自较重原子的反常散射。
\end{solution}

\begin{solution}{pb:b3:x-ray-diffraction:1}
\textbf{1.} $\theta = \ang{14.171}$, $d = \lambda/2\sin\theta = \qty{314.64}{pm}$.
\textbf{2.} 314.64, 222.49, 181.66, 157.32, 140.71, \qty{128.45}{pm}.
\textbf{3.} 1.000, 2.000, 3.000, 4.000, 5.000, 6.000.
\textbf{4.} 1, 2, 3, 4, 5, 6.
\textbf{5.} 是，对应简单立方晶胞的 (100)、(110)、(111)、(200)、(210)、(211)，
其棱长 $a' = \qty{314.6}{pm}$。
\textbf{6.} $N = 7$（以及 15、23\dots），它不能表示为三个平方数之和；第一个
差异将出现在第八条线上。六条线无法区分棱长为 $a'$ 的 P 晶胞
与棱长为 $2a'$ 的 F 晶胞。
\textbf{7.} $N = 4, 8, 12, 16, 20, 24$: (200), (220), (222), (400), (420), (422).
\textbf{8.} $a = 2d_{200} = \qty{629.3}{pm}$.
\textbf{9.} \ce{KCl} (\qty{629.29}{pm}).
\textbf{10.} $F = 4[f(\ce{K+}) - f(\ce{Cl-})]$ 适用于指标全为奇数的情形，两种离子各有 18
个电子，故振幅相互抵消。
\textbf{11.} 两者均可：\ce{Na+}（10）与 \ce{Cl-}（18），\ce{K+}（18）与 \ce{Br-}（36）
散射方式不同；\ce{NaCl} 的 (111) 峰应位于 \ang{27.36}，\ce{KBr} 的峰位于 \ang{23.38}，
二者均在所记录的范围之内。
\textbf{12.} X 射线将 \ce{K+} 与 \ce{Cl-} 视为几乎相同；
在岩盐型排列中相同的离子构成棱长为 $a/2$ 的简单立方阵列，即一个表观晶胞。
\textbf{13.} 四个 \ce{KCl}。
\textbf{14.} $\rho = 4 \times 74.6/(6.022\times10^{23} \times (6.293\times10^{-8})^3) =
\qty{1.99}{g/cm^3}$.
\textbf{15.} 六个与六个（另一种离子形成的八面体）。
\textbf{16.} $a/2 = \qty{314.6}{pm}$.
\textbf{17.} (440)，$N = 32$，位于 \ang{87.65}（全奇指标的 (333)/(511) 缺失）。
\textbf{18.} 强度会随晶粒的择优取向、
吸收、散射因子的衰减以及热运动而变化；
而位置仅取决于点阵。
\textbf{19.} 由 $d = \lambda/2\sin\theta$，$\Delta d/d = -\cot\theta\,\Delta\theta$：相同的角度误差
在较大 $\theta$ 处给出较小的相对误差。
\textbf{20.} $d_{420} = \lambda/2\sin\ang{33.190} = \qty{140.71}{pm}$, $a = d\sqrt{20} =
\qty{629.3}{pm}$.
\textbf{21.} $L = 0.9 \times 0.15406/(4.36\times10^{-3} \times \cos\ang{33.19}) = \qty{38}{nm}$.
\textbf{22.} 否：$\beta = \ang{0.0095}$，远低于典型的仪器宽度。
\textbf{23.} \textbf{$a = \qty{629.3}{pm}$：粉末为氯化钾}，其
奇数指标的反射因两种电子数相同的离子而消光。
\end{solution}
